\section{\appendixHtitle}\label{app:doubling}

\subsection{What is formalised}\label{sec:doubling_formal}

This appendix is an account of a Lean~4 development, checked against Mathlib
\texttt{v4.31.0}. Every statement below is a theorem of that development; every proof sketch
names the results that the corresponding proof term actually invokes, and nothing else. The
development contains no unproved goal, and each result stated here has been verified to depend
on no axiom beyond Lean's own \texttt{propext}, \texttt{Classical.choice} and
\texttt{Quot.sound}.

\medskip\noindent
\emph{The object.} The graph is \extref{the counter-example graph of the main text}: a source
$s_0$, a ladder $1,2,3,\dots$, a sink $s_f$, and the forward edges $s_0\rightarrow1$,
$j\rightarrow j+1$, $2j\rightarrow j$, $j\rightarrow s_f$. The predecessors of a ladder state
$j$ are $j-1$ and $2j$, so the backward chain either \emph{decrements} or \emph{doubles}, and
the loop closure adds the wrap $s_f\rightarrow s_0$. The backward policy is
$\BackwardPolicy(j\rightarrow2j)=\varepsilon(j)$ and
$\BackwardPolicy(j\rightarrow j-1)=1-\varepsilon(j)$, and the family studied throughout is
\begin{equation}\label{eq:doubling_family}
    \varepsilon_{c,s}(j)\;:=\;\frac{c}{(j+1)^{s}},\qquad s\geq0,\quad 0<c<2^{s} .
\end{equation}
The target row $\widehat\BackwardPolicy(s_f\rightarrow\cdot)$ is a probability supported in
$\{1,\dots,d\}$ of mean $\bar\jmath$. A \emph{truncation at $K$} deletes every doubling edge
$j\rightarrow2j$ with $2j>K$, leaving a chain on $K+2$ states.

\medskip\noindent
\emph{What is not modelled.} The development builds no stochastic process. It works with the
invariant measure through the balance identity across a cut, with the expected hitting time
through the supersolution that dominates it, and with the descent through a recursion on the
state rather than a chain of random variables. Everything below is therefore a statement about
those objects, and \S\ref{sec:doubling_differences} says where that costs something.


\subsection{The results}\label{sec:doubling_results}


\subsubsection{The chain, and when it settles}


\begin{theorem}
  \label{thm:main_irreducible}
On the loop closure every state reaches every other, and every invariant measure of it is
strictly positive at every state.
\end{theorem}

\begin{proof}
The graph half is a pair of explicit paths \leanref{reach\_all\_none}. Every state
reaches the ladder state $1$: from a ladder state the decrements walk down to the source, the
wrap carries it to the sink, and the target row re-enters the ladder. Conversely $1$ reaches
every state, by climbing: from $n$ the doubling edge to $2n$ followed by $n-1$ decrements lands
on $n+1$, and $2n\geq n+1$. Concatenating the two gives reachability between any pair.

The measure half is general and needs only that \leanref{PreStat.pos\_of\_irreducible}:
an invariant measure is positive somewhere, and invariance transports that positivity along any
path, so irreducibility spreads it over the whole state space.
\end{proof}

\noindent Irreducibility and positivity are the standing hypotheses under which the rest of this appendix
speaks of \emph{the} invariant probability and of the operators built from it. The main text's
universality and convergence theorems are stated for chains carrying such a measure; what
follows establishes, for each member of the family, whether one exists.


\begin{theorem}
  \label{thm:main_phase}
Let $\varepsilon_{c,s}(j)=c(j+1)^{-s}$ with $c>0$, in the standing range. The loop-closed
backward chain admits an invariant probability exactly in two of the four regimes: it admits
none for $s<1$; exactly one for $s=1$ with $c<1$; none for $s=1$ with $c\geq1$; and one for
$s>1$.
\end{theorem}

\begin{proof}
Three cases, each a drift computation against a Lyapunov function, and no more.

For $s<1$ the drift coefficient $c(m+1)^{1-s}$ grows without bound, which forces the tail of
any candidate invariant sequence to grow too fast to be summable
\leanref{isEmpty\_stat\_of\_family\_lt\_one}. For $s=1$ the coefficient is the constant
$c$, and $c\geq1$ is exactly the threshold past which the same tail argument applies
\leanref{isEmpty\_stat\_of\_family\_ge\_one}. For $s>1$ the drift is eventually
negative and a Foster argument on the ladder height produces an invariant probability
\leanref{exists\_stat\_of\_family\_gt\_one}; the remaining regime $s=1$, $c<1$ is the
same argument at the explicit threshold.

Each conclusion is drawn from the cut balance directly, without constructing the chain as a
stochastic process.
\end{proof}

\noindent This is the phase diagram that decides which members of the family the main text's theorems can
even be stated about. In the two regimes with no invariant probability there is no $\lambda$,
hence no $\Pi$, no mixing coefficients and no diffusion constant: the hypotheses of those
theorems fail there for want of the objects they quantify over, which is a different kind of
failure from the one this appendix is about.


\subsubsection{The invariant measure}


\begin{theorem}
  \label{thm:main_invariant_measure}
At $s=1$ and $0<c<1$, let $\lambda$ be an invariant probability of the loop closure. Then the
Cram\'er equation $\psi_c(p)=0$ has a root $p>1$, and there is a constant $C>0$ with
\[ \lambda_m\,m^{p}\longrightarrow C, \qquad (p-1)\,m^{p-1}L(m)\longrightarrow C, \]
where $L(m)=\sum_{j\geq m}\lambda_j$; and the first convergence carries a polynomial rate:
there are $\vartheta\in(0,1)$ and $c_6>0$ with $|\lambda_m m^{p}-C|\leq c_6\,m^{-\vartheta}$ at
every $m\geq1$.
\end{theorem}

\begin{proof}
The Cram\'er root and the parameters it fixes are packaged by \leanref{Decay.ofC}.
Given them, the invariant probability yields a positive sequence satisfying the cut balance
above $d$, and the whole asymptotic follows from that alone
\leanref{Decay.sharp\_of\_cutBal}: a contraction estimate on dyadic blocks shows the
rescaled profile $\lambda_m m^{p}$ is Cauchy, its limit lies between the two constants of the
two-sided decay bound and is therefore positive, and balancing the block index against the
level gives the exponent $\vartheta$. The statement about the tail is the same limit summed
\leanref{Decay.tendsto\_tail\_sharp}: comparing $\sum_{j\geq m}j^{-p}$ with its
integral turns a limit for the profile into one for the tail.
\end{proof}

\noindent The invariant measure of this graph decays like a power, not geometrically, and the exponent is
the Cram\'er root of the doubling probability. That is the mechanism behind everything that
follows: a power tail is heavy enough that the centred tail indicators below have vanishing
Rayleigh quotient, which is what breaks the mixing hypothesis of the main text.


\begin{theorem}
  \label{thm:main_constant_determined}
At $s=1$ and $0<c<1$, two invariant probabilities of the loop closure that agree on the first
$d$ ladder states agree at every ladder state.
\end{theorem}

\begin{proof}
At $s=1$ the cut balance expresses $\lambda_m$, for every $m>d$, as a positive combination of
the $\lambda_j$ with $\lceil m/2\rceil\leq j\leq m-1$ — indices all smaller than $m$. Induction
on $m$ therefore determines the whole sequence from $\lambda_1,\dots,\lambda_d$, which is the
uniqueness statement \leanref{Decay.cutBal\_unique}; the Cramér root that the
surrounding parameters need is supplied by \leanref{Decay.ofC}.
\end{proof}

\noindent The boundary data of the target row therefore fix the whole invariant measure, and with it the
constant of the asymptotic below. This is what makes that constant a property of the policy and
the target row rather than a free parameter.


\begin{theorem}
  \label{thm:main_constant_functional}
At $s=1$ and $0<c<1$ there is a root $p>1$ of the Cramér equation $\psi_c(p)=0$ and there are
coefficients $\nu_1,\dots,\nu_d\geq0$ with $\sum_{j\leq d}\nu_j>0$ and $\nu_j=0$ whenever
$2j\leq d$, such that for \emph{every} invariant probability $\lambda$ of the loop closure,
\[ \lambda_m\,m^{p}\;\longrightarrow\;\sum_{j=1}^{d}\nu_j\,\lambda_j . \]
\end{theorem}

\begin{proof}
The sequences satisfying the cut balance above $d$ form a real vector space, and restriction to
the first $d$ coordinates is an isomorphism onto $\mathbb R^{d}$ — that is the content of the
determinacy above. Each map $\lambda\mapsto\lambda_m m^{p}$ is linear in those coordinates, so
their limit is additive and positively homogeneous on the cone of positive sequences and
extends to a linear functional on the whole space; its coefficients are the $\nu_j$
\leanref{Decay.constant\_functional}. Positivity of $\sum_j\nu_j$ comes from the
two-sided decay bound, and $\nu_j=0$ for $2j\leq d$ because no window above $d$ ever reads
those indices. The root and the parameters it determines are again
\leanref{Decay.ofC}.
\end{proof}

\noindent The constant of the tail asymptotic is thus a fixed non-negative linear functional of the
boundary data, the same functional for every invariant probability. The vanishing of $\nu_j$
for $2j\leq d$ says that the low half of the target row contributes nothing: only the states
the doubling edge can reach from above carry weight.


\subsubsection{The backward length}


\begin{theorem}
  \label{thm:main_sigmaBar_eq}
At $s=1$ and $0<c<1$ the expected backward-trajectory length of the loop closure is exactly
\[ \bar\sigma \;=\; \frac{\bar\jmath}{1-c}, \]
the supremum of the truncated expectations $\bar\sigma^{(n)}$.
\end{theorem}

\begin{proof}
The linear function $h(x)=x/(1-c)$ is an exact solution of the one-step equation at $s=1$: the
substitution $\varepsilon(j)=c/(j+1)$ collapses the average of the state after one step to
$c+j-1$, and $1+h(c+j-1)=h(j)$. A non-negative solution of that equation dominates the
truncated hitting-time expectations by induction, and the constant drift $c-1$ makes the
inequality an equality in the limit \leanref{sbar\_iSup\_eq}. Averaging over the
target row, whose mean is $\bar\jmath$ and on which $h$ is linear, gives the stated value.
\end{proof}

\noindent This is the quantity Morozov et al.\ (2025) require to be finite, and on this graph it is
finite throughout the positive recurrent regime at $s=1$. It is the first half of the
separation this appendix is for: the condition of the main text that is repaired here.


\begin{theorem}
  \label{thm:main_sigmaBar_le}
At $s=1$ and $0<c<1$, on the loop closure and on every truncation alike, the truncated expected
backward-trajectory lengths satisfy $\bar\sigma^{(n)}\leq\bar\jmath/(1-c)$ for every $n$.
\end{theorem}

\begin{proof}
The same linear supersolution as above serves here, and the one place the truncation differs is
harmless: at a state whose doubling edge has been cut the decrement carries all the mass, and
the one-step equation becomes an inequality in the right direction, because $1\leq1/(1-c)$. The
induction and the target-row average then run unchanged \leanref{sigmaBar\_le}.
\end{proof}

\noindent The bound is uniform in the truncation, so the finiteness that repairs the condition of the
main text is not an artefact of the infinite chain: it survives every finite approximation of
it, which is what the measurements of a truncated system can be compared against.


\subsubsection{The diffusion operator}


\begin{theorem}
  \label{thm:main_unbounded}
Let $\varepsilon_{c,s}$ be a member of the family with $c>0$, $s\geq0$, let $\lambda$ be an
invariant probability of the loop closure, and let $1\leq p<\infty$. Then:
\begin{itemize}
  \item for every $t>0$ there is a cut $m>d$ whose centred tail indicator $f_m$ satisfies
        $\|(\mathrm{Id}-P_\star)f_m\|^p_{L^p(\lambda)}\leq t\,\|f_m\|^p_{L^p(\lambda)}$; and
  \item no constant $B$ satisfies $\|f\|_{L^p(\lambda)}\leq B\,\|(\mathrm{Id}-P_\star)f\|_{L^p(\lambda)}$
        for every bounded mean-zero $f$.
\end{itemize}
\end{theorem}

\begin{proof}
The family satisfies the growth condition $\sum_{i<D}\log(1/\varepsilon(2^i))=o(2^{D})$
\leanref{GrowthCond}, because for $\varepsilon_{c,s}$ that sum is $O(D^2)$
\leanref{growthCond\_of\_family}.

Given the growth condition, the ratio $L(m)/\lambda_m$ is unbounded: were it bounded by $Q$,
the tail would decay geometrically at rate $e^{-1/Q}$, while iterating the doubling inequality
along the powers of two forces $L(2^{D})\geq\lambda_1\prod_{i<D}\varepsilon(2^{i})$, and the
growth condition makes the second bound incompatible with the first. Choosing cuts along that
unbounded ratio makes the centred tail indicator's defect small against its own mass
\leanref{Stat.exists\_small\_mass\_ratio}, which is the first clause; the second is
that same infimum read as the non-existence of a bound
\leanref{Stat.no\_bounded\_inverse}.
\end{proof}

\noindent The second clause is the failure of a bounded diffusion operator at every finite exponent. The
main text's convergence bounds are stated with a constant of exactly this shape; on this graph
no such constant exists, at any $p$, at any positive recurrent member of the family. The first
clause is the sharper fact behind it: the infimum of the Rayleigh quotient is not merely
unattained but zero.


\begin{theorem}
  \label{thm:main_unbounded_infty}
Let $\varepsilon_{c,s}$ be a member of the family with $c>0$ and let $\lambda$ be an invariant
probability of the loop closure. There is no constant $B$ such that every bounded mean-zero
$f$, attaining size $M$ somewhere and having flow-matching defect bounded uniformly by $D$,
satisfies $M\leq B\,D$.
\end{theorem}

\begin{proof}
The witnesses are clipped ramps: $r_N$ is the ladder height cut off at $N$. Its
flow-matching defect is bounded independently of $N$ — at a state below the clip the defect is
$1-\varepsilon(j)(j+1)$, at a state above it the defect vanishes, and at the sink it is
$-\bar\jmath$ — while centring leaves the defect unchanged and leaves the function attaining at
least $N/2$. So $M/D$ grows without bound
\leanref{Stat.no\_bounded\_inverse\_infty}.

That the defect is bounded uses $\gamma=\sup_j(j+1)\varepsilon(j)<\infty$, which for this
family holds exactly when $s\geq1$, with $\gamma=c\,2^{1-s}$
\leanref{gamma\_family}. Below $s=1$ there is no invariant probability to quantify
over \leanref{isEmpty\_stat\_of\_family\_lt\_one}, and the statement holds for that
reason instead.
\end{proof}

\noindent The failure therefore reaches the uniform norm as well, where the main text's stable bounds are
most often read. No hypothesis on $s$ is needed: above $s=1$ the clipped ramps do the work, and
below it there is no invariant probability at all, so the claim is vacuous exactly where the
objects it quantifies over do not exist.


\begin{theorem}
  \label{thm:main_unbounded_two}
Let $\varepsilon_{c,s}$ be a member of the family with $c>0$, $s\geq0$, and let $\lambda$ be an
invariant probability of the loop closure. On $L^2(\lambda)$: no bounded operator $S$ satisfies
$S(\mathrm{Id}-P_\star)=\mathrm{Id}-\Pi$; the series $\sum_{n\geq0}(P_\star^n-\Pi)$ does not
converge in operator norm; and $\sum_{n\geq0}\widehat\beta_n=+\infty$, where
$\widehat\beta_n=\|P^n-\Pi\|_{L^2(\lambda)}$.
\end{theorem}

\begin{proof}
All three are the same fact at $p=2$. The family satisfies the growth condition
\leanref{GrowthCond}, discharged as before
\leanref{growthCond\_of\_family}; the vanishing infimum of the Rayleigh quotient
then excludes a bounded left inverse \leanref{Stat.no\_diffusionOp}. The other two
follow from that one: a norm-convergent series would furnish such an inverse, since its sum
$U$ satisfies $(\mathrm{Id}-P_\star)U=U(\mathrm{Id}-P_\star)=\mathrm{Id}-\Pi$ by telescoping
\leanref{Stat.not\_tendsto\_partialSum}; and a summable
$\sum_n\widehat\beta_n$ would make that series absolutely convergent, the norms being the
$\widehat\beta_n$ \leanref{Stat.not\_summable\_betaHat}.
\end{proof}

\noindent This is the sentence that separates the two conditions. The summable-mixing hypothesis of the
main text's $L^2$ universality theorem, and of its frozen detailed-balance theorem in the
flow-matching instance --- where the constant is exactly the state-space mixing sum
$\sum_n\widehat\beta_n$ --- is repaired by no positive recurrent member of this family, while
the backward length above is finite throughout. The divergence is a statement about the
flow-matching instance; the detailed-balance instance, whose coefficients are the edge-lifted
ones, is not covered.


\begin{theorem}
  \label{thm:main_rate}
At $s=1$ and $0<c<1$, with $p>1$ the Cram\'er root and $\lambda$ an invariant probability of
the loop closure, there are constants $0<c_1\leq c_2$ and a cut $m_2>d$ such that the centred
tail indicators satisfy
\[ c_1\,m\;\bigl\|(\mathrm{Id}-P_\star)f_m\bigr\|^2_{L^2(\lambda)}
   \;\leq\; 4c_2(p-1)\,\|f_m\|^2_{L^2(\lambda)} \qquad (m\geq m_2). \]
\end{theorem}

\begin{proof}
The two-sided decay bound $c_1 j^{-p}\leq\lambda_j\leq c_2 j^{-p}$ holds unconditionally at
this level \leanref{Decay.decay\_two\_sided\_of\_cutBal}, run at the level $m_0$ at
which the descent-weight window is controlled \leanref{Decay.m0}, with the root
supplied by \leanref{Decay.ofC}. Summing that bound over $j\geq m$ gives
$L(m)/\lambda_m\geq c_1m/(c_2(p-1))$, so the per-cut inequality --- the defect of $f_m$ is
carried by the cut and the window alone, and the cut balance collapses the window --- reads as
the display once the mass below the cut is bounded away from zero. The cut $m_2$ past which
that last condition holds exists because $L(m)\to0$
\leanref{Stat.exists\_rayleigh\_family}.
\end{proof}

\noindent Unboundedness above says only that the Rayleigh quotient has infimum zero; this says how fast,
along an explicit family of trial functions. Squared and rearranged it is
$\|f_m\|\geq c_7\sqrt m\,\|(\mathrm{Id}-P_\star)f_m\|$. It is not a lower bound on the
diffusion constant of the infinite chain, which is $+\infty$ there; it is a rate along the
$f_m$, and it is the truncation below that converts a rate into a bound on a finite constant.


\subsubsection{Flow matching}


\begin{theorem}
  \label{thm:main_unsolvable}
Let $\varepsilon_{c,s}$ be a member of the family with $c>0$, $s\geq0$, and let $\lambda$ be an
invariant probability of the loop closure. Then $(\mathrm{Id}-P_\star)L^2(\lambda)$ is a dense
proper subspace of $\ker\Pi$, and there are probability densities
$\finit,\fterm\in L^2(\lambda)$ for which no $f\in L^2(\lambda)$ satisfies the flow-matching
equation $(\mathrm{Id}-P_\star)f=\finit-\fterm$.
\end{theorem}

\begin{proof}
Write $A$ for the restriction of $\mathrm{Id}-P_\star$ to $\ker\Pi$. It is injective, its
kernel being the mean-zero fixed points of $P_\star$, which irreducibility reduces to the
constants and hence to zero; its adjoint inside $\ker\Pi$ is the restriction of
$\mathrm{Id}-P$, injective for the same reason, so the range of $A$ is dense. It is not
surjective: a surjective bounded injection of a Hilbert space is boundedly invertible by the
open mapping theorem, and that contradicts the vanishing infimum established above. Picking
$\psi\in\ker\Pi$ outside the range and splitting it into positive and negative parts produces
the two densities \leanref{doubling\_unsolvable}. The growth condition the argument
needs is discharged for the family as before
\leanref{growthCond\_of\_family}.
\end{proof}

\noindent What this refutes is \emph{exact} $L^2$ flow matching on the infinite chain: a pair of $L^2$
densities the loop closure cannot balance at all. It does not refute the weak universality of
the main text, whose conclusion is a defect infimum of zero --- density of the range is exactly
what the first clause asserts, and whether the weak form fails here is not decided.


\subsubsection{The truncation}


\begin{theorem}
  \label{thm:main_truncation_irreducible}
At an even cap $K\geq d$ the truncation is irreducible on its $K+2$ states, and every invariant
measure of it is strictly positive at each of them.
\end{theorem}

\begin{proof}
Only the climbing half differs from the infinite chain, because the doubling edges above $K/2$
have been cut. Let $2^{r}$ be the largest power of two at most $K$: the edges
$2^{i}\rightarrow2^{i+1}$ all survive for $i<r$, so the ladder state $1$ reaches $2^{r}$ and,
by decrements, every state below it. Above $2^{r}$ an even $n$ is reached by doubling from
$n/2\leq K/2<2^{r}$, and an odd $n$ is not the cap --- which is even --- so $n+1$ is even and
in range and one decrement finishes \leanref{reach\_all\_some}. Positivity is the
same general fact as on the loop closure
\leanref{PreStat.pos\_of\_irreducible}.
\end{proof}

\noindent The truncation is what a numerical experiment can actually run, and this is what makes it a
legitimate finite approximation: the same qualitative structure as the infinite chain, on a
finite state space where the diffusion constant is finite and can be measured.


\begin{theorem}
  \label{thm:main_truncation_bhat}
At an even cap $K\geq d$ the truncation carries exactly one invariant probability $\lambda^K$,
and on $L^2(\lambda^K)$ there is a bounded operator $S$ with
\[ S(\mathrm{Id}-P_\star)=(\mathrm{Id}-P_\star)S=\mathrm{Id}-\Pi,\qquad \Pi S=S\Pi=0 . \]
In particular the diffusion constant $\widehat B_K=\|S\|$ is finite at every even cap.
\end{theorem}

\begin{proof}
Existence and uniqueness of the invariant probability are those of a finite irreducible chain
\leanref{exists\_stat}, \leanref{stat\_unique}. On a finite state space
$\mathrm{Id}-P_\star+\Pi$ is then injective --- a vector it kills has zero mean and is fixed by
$P_\star$, hence is a constant of zero mean --- and injective is invertible in finite
dimension; writing $R$ for the inverse and $S=R-\Pi$, the relations $\Pi P_\star=P_\star\Pi=\Pi$
and $\Pi^2=\Pi$ give the four identities \leanref{Stat.exists\_bhat}. Boundedness
is automatic, the space being finite-dimensional.
\end{proof}

\noindent So the object that fails to exist on the infinite chain does exist on every truncation of it.
Nothing diverges at a finite cap; what the next result shows is that the constant grows without
bound as the cap does, which is how an infinite constant is visible to a finite computation.


\begin{theorem}
  \label{thm:main_truncation_sqrtK}
At $s=1$ and $0<c<1$ there are $c_8>0$ and a cap $K_0$ such that, for every $K\geq K_0$ and
every bounded operator $S$ on $L^2(\lambda^K)$ with $S(\mathrm{Id}-P_\star)=\mathrm{Id}-\Pi$,
\[ c_8\,K \;\leq\; \|S\|^2 . \]
That is, $\widehat B_K\geq\sqrt{c_8}\,\sqrt K$ at every large even cap.
\end{theorem}

\begin{proof}
The argument is the per-cut inequality applied at a cut of order $K$. The decay bound transfers
to the truncation with constants that do not depend on the cap
\leanref{Stat.decayK}: extending $\lambda^K$ above the cap gives a positive
sequence satisfying the cut balance below $K/2$, so the block estimate applies at the level
$m_0$ \leanref{Decay.m0}, and the ratio of the two constants is controlled by
chaining the cut-balance inequalities across one block. The descent weight at that level is
two-sided \leanref{Decay.descOne\_two\_sided} with the escape constant
\leanref{Decay.c5}, and the Cram\'er root is again
\leanref{Decay.ofC}.

Taking the cut at $\lfloor K/4\rfloor$ then makes $L_K(m)/\lambda^K_m$ of order $K$, while the
mass below the cut is bounded away from zero uniformly in $K$ --- on a finite irreducible chain
the mass at the source is the reciprocal of the expected return time, which is
$2+\bar\sigma_K\leq2+\bar\jmath/(1-c)$. Feeding both into the per-cut inequality gives the
displayed bound \leanref{Stat.sqrtK}.
\end{proof}

\noindent This is the quantitative form of the separation, and the one a measurement can be compared
against: the diffusion constant of the truncation is finite but grows at least like $\sqrt K$,
so it diverges as the truncation is lifted, while the backward length stays at
$\bar\jmath/(1-c)$ throughout. The bound holds for \emph{every} left inverse, so it is a
property of the chain and not of a particular construction.


\subsection{What is not established here}\label{sec:doubling_differences}

Three claims that would naturally accompany the results above are absent, and in each case for
the same reason: they are properties of the backward chain as a stochastic process, and no such
process is constructed.

\begin{itemize}
  \item \emph{Transience and null recurrence.} The phase diagram distinguishes only between the
        existence and the non-existence of an invariant probability. Where none exists --- at
        $s<1$, and at $s=1$ with $c\geq1$ --- nothing is said about which of transience or null
        recurrence obtains. In particular the case $s=1$, $c=1/\ln2$ sits inside the
        non-existence regime and is not resolved further.
  \item \emph{An infinite expected hitting time.} At $s=1$ with $1\leq c<2$ the absence of an
        invariant probability is established; the accompanying statement that
        $\mathbb E(\sigma\mid X_0=j)=+\infty$ at every ladder state is not.
  \item \emph{The trajectory reading of the diffusion operator.} The operator excluded at $p=2$
        is characterised by the resolvent identities, not as the sum
        $\sum_{n\geq0}\bigl(\mathbb E(\theta(X_n)\mid X_0=x)-\lambda(\theta)\bigr)$ along
        backward trajectories. That the two agree needs the chain.
\end{itemize}

\noindent
Nothing above depends on these. They are named so that a reader who expects the vocabulary of
recurrence classification knows where it has been replaced by the weaker statement that is
actually proved.


\subsection{The development}\label{sec:doubling_development}

Each result above is a declaration of a Lean~4 development checked against Mathlib
\texttt{v4.31.0}. Table~\ref{tab:doubling_decls} gives the correspondence. Every one of them
has been verified to rest on no axiom beyond Lean's own \texttt{propext},
\texttt{Classical.choice} and \texttt{Quot.sound}, and the development contains no unproved
goal.

\begin{table}[htb]
\centering
\small
\begin{tabular}{@{}l l r r@{}}
\hline
result & declaration & named results & declarations \\
\hline
\ref{thm:main_irreducible} & \texttt{\small main\_irreducible} & 3 & 67 \\
\ref{thm:main_phase} & \texttt{\small main\_phase} & 25 & 162 \\
\ref{thm:main_invariant_measure} & \texttt{\small main\_invariant\_measure} & 98 & 366 \\
\ref{thm:main_constant_determined} & \texttt{\small main\_constant\_determined} & 5 & 98 \\
\ref{thm:main_constant_functional} & \texttt{\small main\_constant\_functional} & 104 & 377 \\
\ref{thm:main_sigmaBar_eq} & \texttt{\small main\_sigmaBar\_eq} & 10 & 83 \\
\ref{thm:main_sigmaBar_le} & \texttt{\small main\_sigmaBar\_le} & 4 & 39 \\
\ref{thm:main_unbounded} & \texttt{\small main\_unbounded} & 16 & 128 \\
\ref{thm:main_unbounded_infty} & \texttt{\small main\_unbounded\_infty} & 11 & 116 \\
\ref{thm:main_unbounded_two} & \texttt{\small main\_unbounded\_two} & 60 & 267 \\
\ref{thm:main_rate} & \texttt{\small main\_rate} & 71 & 302 \\
\ref{thm:main_unsolvable} & \texttt{\small main\_unsolvable} & 44 & 203 \\
\ref{thm:main_truncation_irreducible} & \texttt{\small main\_truncation\_irreducible} & 4 & 69 \\
\ref{thm:main_truncation_bhat} & \texttt{\small main\_truncation\_bhat} & 37 & 183 \\
\ref{thm:main_truncation_sqrtK} & \texttt{\small main\_truncation\_sqrtK} & 96 & 393 \\
\hline
\end{tabular}
\caption{The results of this appendix as declarations of the Lean development, in the
namespace \texttt{GFNBounds.Doubling}. The last two columns count what each proof rests on,
unfolded to the bottom of the development: the declarations it reaches, and how many of those
are results the development names rather than steps internal to one.}
\label{tab:doubling_decls}
\end{table}

